The detecting value must have a depth bound common to all
primitive source combinations. We choose this bound from the
weight data before choosing a detecting polynomial. Lemma~\ref{lem:lift-plane-symbol}
then lifts the resulting flat symbol to the integral polynomial lattice.
\paragraph{A depth fixed before the choice of a source.}
\label{par:plane-uniform}
\hypertarget{source-plane-return}{}
Let
\[
 L=Nk^2\log_p a.
\]
Here \(a\) denotes the lower bound for
the flat degrees on the tops, as in Section~\ref{par:plane-rank}.
By the inequality on the selected parameters \(a\ge\max(p,4N,k,2)\), and the full
degrees are at most \(2a\). There are at most \(N\) true letters
in the signed model: their positions in the cycle are distinct tops.
Each quotient degree is at most \(2a\) divided by its low radix.
These statements hold before the digits, projections or unit
monomial of the source have been chosen.

\begin{proposition}[An integral lift of the detecting symbol at uniform depth and order]\label{prop:plane-uniform-symbol-return}
Under the hypotheses on the frames and the parameter sizes, let \(X\) be the normal degree bound in the proof and put
\[
 D_X=\lceil49Nk^2\log_p a\rceil,\qquad t_F=\lfloor X\rfloor+Nc_0.
\]
 For each primitive polynomial
in the exact model span there is a detecting flat symbol
of order at most \(t_F\) at depth \(D_X\), realizable
at the selected plane with competitor conductor \(c_P\).
Both \(D_X\) and \(t_F\) are fixed by the weight and size
data before choosing the primitive polynomial.
\end{proposition}

\begin{proof}
\proofstep{The common bounds before choosing a source.}
In \eqref{eq:plane-slice-depth} take for \(D_{\rm no}\) the sum of
the integral projector exponents. At a complement cut of total
degree \(D\), the trace projector bound gives
\(\binom K2\lfloor\log_p(2D+K-1)\rfloor\), interpreted as zero
for a scalar action. The estimate also holds in the factorial
norm by self-adjointness. There are at most \(2N\) such cuts,
since the segment with \(h+r\) listed tops has \(h+r+1\) cuts.
Now
\[
 2D+K\le4Na+k\le a^3;
\]
hence \(D_{\rm no}\le6L\).
Counting carries at each true letter bounds the plane and
complement factorial deficiencies by \(3Nr\log_p a\) and
\(3NK\log_p a\), respectively. Their halves and the
projector loss thus give \(D_n\le9L\).
The full factorial loss \(\kappa_{\rm fl}\) is at most
\(Nk\log_p(2a)\le 2Nk\log_p a\le20L\).

Similarly \(e_d\le4k^2\log_p a\) and
\(K_d\le4Nk\log_p a\), since \(2\sum_i d_i+r\le a^3\).
There are at most \(N\) links and at most \(N\) chain cuts,
so \[
 D_b=D_m+D_l+3K_d\le20L.
\] We may take the integer
\[
                         D_X=\lceil49L\rceil\le50L .
\]
Choose the absolute constant \(C\) in \(D_p\ge CL\) before
taking the host sufficiently large. The conductor hypothesis gives
\[
 c_P\ge c_\star\ge8(D_p+1)>\max\{2D_X,4D_b\},
 \qquad D_X\le50L,\quad D_b\le20L .
\]

For the second reserve in \eqref{eq:plane-gram-reserves}, write
\(D=z_j-a_{j,1}\) and \(P=p^{z_{j+1}-z_j}\). Intervening letters
have distinct distances \(m\ge D+L_1+1\) from \(z_j\).
On a chain cut in segment \(j+1\) the oldest even predecessor
lies at or after \(a_{j,1}\). These two facts give
\[
 b_j\le 2a\sum_{m\ge D+L_1+1}p^{-m}
     \le 3a p^{-D-L_1-1},
 \qquad B_u\ge \frac{a}{p^D P}.
\]
Set \(x=B_u/(Yk^2)\), so \(s_u=\lfloor x\rfloor\).

At the tested letters \(B_u\ge a/Q\) and the host size bound
give \(s_u\ge100r^2 c_P\); in particular \(x\le2s_u\).
The spacing bound \(p^{L_1+1}\ge48rYk^2\) follows from the
host spacing with \(r\le Y\) and \(N\ge1\). Hence
\begin{equation}\label{eq:plane-cut-budget}
 \frac{rb_j}{P}\le \frac{3rB_u}{p^{L_1+1}}
       \le\frac{x}{16}\le\frac{s_u}{8},\qquad
 26r c_P\le\frac{s_u}{4},\qquad 128c_P\le s_u .
\end{equation}
The last two inequalities use \(r\ge2\). Adding the first two costs in \eqref{eq:plane-cut-budget} gives
\[
 rb_j/P+26r c_P\le3s_u/8
\]
on the same cut.

\proofstep{An integral symbol from the model polynomial.}
\label{par:plane-reassembly}
Use the integral model polynomial \(Z_*\) given by the Gram
estimate, and the complement test \(N_0\). Give each monomial
in the model the segment positions of its letters. A block letter
has its intervening position, \(L_{jm}\) and \(T_{jm}\) both have
the outgoing position of that shared row, and \(W_{jm}\) has
its incoming position in the next segment. Keep segment indices
on variables even when their positions in the cycle agree. Place
the complement variables only at the intervening and outgoing
positions. These renamings commute with the signed action on a
prefix and its inverse image. This follows on the variables;
the frame rows are all even and translation by \(n'\) preserves
parity. Coincidence of the label and tangent positions causes
no loss for this assertion.

\proofstep{Prefix orders in the model plane.}
For the renamed plane polynomial the prefix before \(x_u\)
has order at least
\[
                        t_u=T_u-s_u/2-60c_P .
\]
For \(u\le h\) the inverse image is the chain cut. The initial lower bound for the order is \(T_u-s_u/8\).
The two losses in the Gram construction are at most \(3s_u/8\) and
\(24c_P\), respectively. The resulting order is at least
\[
 (T_u-s_u/8)-(3s_u/8+24c_P)=T_u-s_u/2-24c_P\ge t_u .
\]

For \(u>h\) the inverse
image consists of the model link and the outgoing label letters
at earlier rows. These extra letters are vectors. Enlarging the
torus action to them multiplies the previous transform by a
commuting substitution with nonnegative powers, so the lower
bound
\[
 T_u-3s_u/8-60c_P
\]
suffices. For the prefix through an
intervening letter the chain cut at \(u+1\) gives \(t_{u+1}\).
The renamed complement polynomial has upper orders at most
\(k_u+4c_P\) and
\[
 k_{u+1}+4c_P
\]
before and through \(x_u\).

These bounds hold on the full polynomials over \(O\). To use them
segment by segment, take homogeneous components in each
segment letter on the plane side. The complement letter degrees
are already fixed. In one component all the permitted coefficients
in a particular segment, after fixing monomials in the other
segments, belong to the same integral submodule given by the
prefix inequalities. In fact taking such coefficients commutes
with the group of the specified segment. This submodule is
saturated in the finite free module of the given degrees:
its inequalities are vanishing conditions on Taylor or Laurent
coefficients over the fraction field. The quotient is torsion free,
hence free over the DVR.

For the fixed segment degrees, let \(M_j\) be the full
polynomial lattice and \(S_j\subset M_j\) the saturated
submodule cut out by its prefix bounds. Choose an integral
retraction
\[
 r_j:M_j\longrightarrow S_j,\qquad r_j|_{S_j}=\mathrm{id}_{S_j},
\]
possible since the quotient
is free. Fixing all monomials outside \(j\) lands in \(S_j\),
so the endomorphism \(r_j\) in factor \(j\) fixes the component
\(A\). These endomorphisms commute and hence
\[
                   A=(\bigotimes_j r_j)A
                           \in\bigotimes_j S_j.
\]
Applying the same map to a monomial expansion produces
a finite integral sum of products of permitted factors.
The identical argument on the complement introduces no
denominators there either.

\proofstep{Extracting factors with the true degrees.}
Write \(A_j,G_j\) for the factors of a resulting pair. Both have
the designated letter degrees. For an even intervening \(x_u\) take \(M=k_{u+1}+4c_P\),
the upper order of \(G_j\) through the letter. The plane
order before it is at least \(t_u\), and
\[
 t_u-M\ \ge\ m_*+s_u/2-64c_P\ \ge\ m_* .
\]

For an odd intervening \(x_u\) take \(M=k_u+4c_P\)
before the letter. The plane order through it is at least
\(t_{u+1}\) and
\(t_{u+1}-M\ge m_*+s_{u+1}/2-64c_P\ge m_*\).
The lower bounds follow from the respective inequalities
for \(T_u,T_{u+1}\).
The absolute lower bound before each letter is
\(B_u/2-s_u\), since \(T_u\ge B_u/2\).
These are exactly the two parity comparisons and absolute
bound used in the determinant extraction.

For clarity, consider the degree conditions in that extraction.
At an intervening position \(y\), let \(d_{\rm low}(y)\) be the
outer digit degree and \(d_A(y),d_G(y)\) the degrees of these two
factors. The class equation gives
\[
  d_{\rm low}(y)+p^{z_j-y}\bigl(d_A(y)+d_G(y)\bigr)\le\deg(y).
\]
Thus the separate order contributions can be paid out of their sum. At an incoming position
the only letter is \(W_{j-1,m}\), of degree \(b\), and the
complementary degree is zero.

Let \(R_{\rm low}\) be the original low radix and \(R_{\rm seg}\) the incoming segment radix at this letter. Substitution first multiplies the exponent by \(P_{j-1,m}\); passing to true degree then gives
\[
 b\ \longmapsto\ P_{j-1,m}b
      \ \longmapsto\ R_{\rm low}P_{j-1,m}b=R_{\rm seg}b .
\]
This is the contribution from the incoming letter alone. Other terms at the shared letter have nonnegative degree, so if \(d_{\rm coll}\) is its total collapsed plane degree and \(d_{\rm true}\) its true degree, then
\[
 R_{\rm seg}b\le R_{\rm low}d_{\rm coll}\le d_{\rm true}.
\] This argument applies
to every model monomial used, each of which has the prescribed
collapsed degrees. After collecting terms and projecting to a
segment multidegree the same upper bound holds. In particular
the upper allowance
\[
 1.01\deg(y)/p^{z_j-y}
\]
holds in both
places. The extraction proof in Section~\ref{sec:plane-lifts}
only uses this upper inequality when estimating the earlier rows.
It therefore applies also to components which have lost much of
their degree in the incoming frame. This is the one-sided form
of the realization lemma.

\proofstep{Recovering the original low coordinates.}
Multiply each product of segment factors by the original low
label after applying the segment power substitutions. The sum
is \(\operatorname{Asm}_\rho(\phi Z_*,N_0)\): the extra power
on an incoming row is exactly \(P_{j-1,m}\) and the label has
zero support off the tops. Within each fixed component the
result has the required full degrees.

In each true letter the substitutions multiply a segment degree by its radix. Summing over the contributing parts \(b\) therefore gives 
\[
 d_{\rm collapsed}=\sum_b P_b d_b,
\]
 where \(P_b,d_b\) are the radix and degree of the intervening, incoming or outgoing part. Restoring the label then gives the original degree
by the class equation. Its normal degree is also constant: at
an even top add its normal low digits to the low radix times the
complement quotient degree; at an odd top use its plane digits
and plane quotient degree. Sum these numbers over the tops.
If \(S\) is the component of order \(o\), the one-sided realization gives an integral \(P\) with
\[
 P|_{\rm chosen}\equiv w^o S\pmod{w^{o+c_P}},\qquad P|_{\rm competitor}\equiv0\pmod{w^{o+c_P}}
\]
at every competing plane. This proves the symbol membership and its realization on the stated one-sided range.

\proofstep{Keeping one order for all tested generators.}
Finally fix an order before choosing \(H\). The counted generators
have normal degree bounded by
\[
 X=N(\mathcal L a/(pk)+10a p^{-L_1}+kQ)
\]
by
Section~\ref{par:plane-rank}. The common factor has normal degree zero. For a counted generator \(A\) and a product of \(j\le Nc_0\) row operators \(D_1,\ldots,D_j\), the per-operator increment bound gives
\[
 \deg_{\rm normal}(D_1\cdots D_j A)
     \le \lfloor X\rfloor+j\le t_F,\qquad t_F=\lfloor X\rfloor+Nc_0,
\]
whenever the result is nonzero. Taking the span preserves this upper bound. The Fock form pairs only equal exponent tuples.

If \(A=\sum_{o\le t_F}A_o\) is a tested input and
\(Z=\sum_o Z_o\) is the assembled symbol, then the diagonal
Fock pairing gives
\[
                  [A,Z]=\sum_{o\le t_F}[A_o,Z_o] .
\]
The discarded terms have disjoint exponent support from \(A\).
After this deletion the ultrametric inequality selects a
summand with pairing valuation no greater than that of the sum. We have proved the flat symbol estimate
at the same \(D_X,t_F,c_P\) for every primitive polynomial in
the exact model span. \label{cor:plane-return-1}
\end{proof}
